\scanpage{181}\mspage{36}%
pouvoir se paraphraser dans le cadre des schémas excellents, à grands coups de voisinages tubulaires étales. Dans l'un et l'autre cas, le ``type d'homotopie fin'' d'un espace modéré, resp. d'un schéma excellent, se définirait par passage à la limite à partir des ``\inserted{$\unsure{\mathcal P|P}$-}types d'homotopie'' associés à des stratifications équisingulières (à strates non sing.) $\mathcal P$ de plus en plus fines. Ce type d'homotopie \inserted{``fin''} incarnerait la connaissance, non des seuls faisceaux ou \inserted{$n$-}champs \unsure{loc.|loc} constants, mais \inserted{(via un passage à la limite \unsure{inductive})} la connaissance de \uline{tous}. Et il dépendrait, en un sens convenable,
% A distinct low inline point after fonctoriellement is preserved literally; its origin and punctuation function are not established (whole-image approximately 700,1365).
fonctoriellement. de $X$. Dans%
% Detached nontext point below the end of Dans, approximately (1170,1369); authorship/origin uncertain. The inline zero-width drawing does not assign it a textual role.
\raisebox{0pt}[0pt][0pt]{\parbox{0pt}{\figure{\begin{tikzpicture}[overlay]\fill (0,-0.75ex) circle[radius=0.35pt];\end{tikzpicture}}}} le cas d'un schéma de type fini sur un corps alg. clos disons, les \uline{théorèmes de finitude} cohomologique ou homotopique les plus forts s'exprimeraient précisément en termes des types d'homotopie fin, et diraient que les types d'homotopie \unsure{ordinaire|ordïnaire} qui en sont les élément constituants sont essentiellement des ``polyédres finis'' -- et $\overline m$ des variétés à bord compactes -- ou de façon plus précise,
% archive sheet number: 181
